\section{项目背景与目标}
\label{sec:introduction}
\subsection{研究背景与任务目标}
Breakout 的动作空间很小，但控制过程同时包含视觉感知、延迟奖励与连续决策。挡板移动通常不会立即得分；一次接球的收益可能在小球多次飞行和碰撞后才体现。智能体需要从连续图像中判断运动趋势，并将后续击砖收益与先前动作联系起来。

本项目在统一的自制固定球速 Breakout 环境中，实现 DQN 和 PPO 两条图像控制链路，研究动作价值学习与直接策略优化的原理、实现和训练表现。任务要求智能体仅根据图像选择四个离散动作，在完整五生命游戏中提高原生得分。环境规则由项目代码定义，结果对应这一任务，不直接对照公开 Atari 基准。

\subsection{研究内容与报告安排}
项目工作包括环境与输入接口建设、两种算法的训练与恢复、开发集模型选择、固定模型测试，以及网页实时展示。DQN 重点研究低学习率续训、训练退化与恢复实验；PPO 重点研究策略与价值网络、更新稳定性及 GAE 单因素对照。实验同时保留有效结果与负结果，以日志、检查点和逐局分数支持分析。

第\ref{sec:shared}节给出共用环境、系统接口和评价口径；第\ref{sec:dqn}节和第\ref{sec:ppo}节分别完整介绍 DQN、PPO 的原理、实现、训练及开发分析；第\ref{sec:comparison}节集中报告冻结模型的统一测试与方法比较；第\ref{sec:conclusion}节总结结论和局限。开发数据用于选模，最终测试用于评价已冻结的候选，二者分别报告。

\FloatBarrier
